\textbf{Phần 1. Trắc nghiệm lựa chọn}
\setcounterpageref{ex}{0}

\begin{ex} %[5H9V49]
Một hình khai triển của hình lập phương gồm $6$ hình vuông bằng nhau, mỗi hình vuông có chu vi là $24\text{ cm}$. Khi gấp lại thành hình lập phương thì diện tích toàn phần của hình lập phương đó là
\choice
{$144\text{ cm}^2$}
{\True $216\text{ cm}^2$}
{$36\text{ cm}^2$}
{$288\text{ cm}^2$}
\loigiai{
Cạnh của hình vuông trong hình khai triển là:
$$24 : 4 = 6\text{ (cm)}.$$
Diện tích của một mặt là:
$$6 \times 6 = 36\text{ (cm}^2\text{)}.$$
Diện tích toàn phần của hình lập phương bằng tổng diện tích $6$ mặt:
$$36 \times 6 = 216\text{ (cm}^2\text{)}.$$
\textbf{Đáp án của câu hỏi là B}
}
\end{ex}

\begin{ex} %[5H9V49]
Khai triển mặt xung quanh của một hình trụ có chiều cao $10\text{ cm}$ ta được một hình chữ nhật có chu vi $48\text{ cm}$. Biết chiều rộng của hình chữ nhật chính là chiều cao hình trụ. Chu vi đáy của hình trụ đó là
\choice
{$28\text{ cm}$}
{\True $14\text{ cm}$}
{$19\text{ cm}$}
{$38\text{ cm}$}
\loigiai{
Nửa chu vi của hình chữ nhật là:
$$48 : 2 = 24\text{ (cm)}.$$
Chiều dài của hình chữ nhật (chính là chu vi đáy hình trụ) là:
$$24 - 10 = 14\text{ (cm)}.$$
\textbf{Đáp án của câu hỏi là B}
}
\end{ex}

\begin{ex} %[5H9V49]
Để làm một chiếc hộp hình lập phương không nắp có cạnh $8\text{ cm}$, bạn Mai cần cắt một tấm bìa thành hình khai triển gồm bao nhiêu hình vuông và tổng diện tích của tấm bìa đó ít nhất là bao nhiêu xăng-ti-mét vuông (không kể mép dán)?
\choice
{$6\text{ hình vuông}, 384\text{ cm}^2$}
{\True $5\text{ hình vuông}, 320\text{ cm}^2$}
{$5\text{ hình vuông}, 256\text{ cm}^2$}
{$4\text{ hình vuông}, 256\text{ cm}^2$}
\loigiai{
Vì hộp hình lập phương không nắp nên hình khai triển gồm $5$ hình vuông.\\
Diện tích một hình vuông là:
$$8 \times 8 = 64\text{ (cm}^2\text{)}.$$
Tổng diện tích bìa tối thiểu là:
$$64 \times 5 = 320\text{ (cm}^2\text{)}.$$
\textbf{Đáp án của câu hỏi là B}
}
\end{ex}

\begin{ex} %[5H9V49]
Hình khai triển của một hình hộp chữ nhật có kích thước dài $12\text{ cm}$, rộng $8\text{ cm}$, cao $5\text{ cm}$. Tổng diện tích của $6$ mặt trên hình khai triển là
\choice
{$384\text{ cm}^2$}
{\True $392\text{ cm}^2$}
{$480\text{ cm}^2$}
{$196\text{ cm}^2$}
\loigiai{
Tổng diện tích $6$ mặt của hình khai triển chính là diện tích toàn phần của hình hộp chữ nhật:\\
Diện tích xung quanh là:
$$(12 + 8) \times 2 \times 5 = 200\text{ (cm}^2\text{)}.$$
Diện tích hai mặt đáy là:
$$12 \times 8 \times 2 = 192\text{ (cm}^2\text{)}.$$
Tổng diện tích là:
$$200 + 192 = 392\text{ (cm}^2\text{)}.$$
\textbf{Đáp án của câu hỏi là B}
}
\end{ex}

\begin{ex} %[5H9V49]
Một hình trụ có bán kính đáy là $3\text{ cm}$ và chiều cao $7\text{ cm}$. Khi khai triển toàn phần của hình trụ, ta được một hình chữ nhật và hai hình tròn. Diện tích của hình chữ nhật biểu diễn mặt xung quanh là (lấy $\pi = 3{,}14$)
\choice
{$65{,}94\text{ cm}^2$}
{\True $131{,}88\text{ cm}^2$}
{$197{,}82\text{ cm}^2$}
{$188{,}4\text{ cm}^2$}
\loigiai{
Chu vi đáy của hình trụ (chiều dài của hình chữ nhật) là:
$$3 \times 2 \times 3{,}14 = 18{,}84\text{ (cm)}.$$
Diện tích của hình chữ nhật xung quanh là:
$$18{,}84 \times 7 = 131{,}88\text{ (cm}^2\text{)}.$$
\textbf{Đáp án của câu hỏi là B}
}
\end{ex}

\begin{ex} %[5H9V49]
Trên các mặt của hình khai triển của một khối xúc xắc hình lập phương được ghi các số từ $1$ đến $6$. Khi gấp lại thành khối xúc xắc chuẩn, tổng số chấm ở hai mặt đối diện luôn bằng $7$. Mặt đối diện với mặt ghi số $3$ phải ghi số nào?
\choice
{$5$}
{\True $4$}
{$2$}
{$6$}
\loigiai{
Theo quy tắc xúc xắc chuẩn, tổng số chấm ở hai mặt đối diện bằng $7$.\\
Do đó, mặt đối diện với mặt ghi số $3$ là:
$$7 - 3 = 4.$$
\textbf{Đáp án của câu hỏi là B}
}
\end{ex}

\begin{ex} %[5H9V49]
Một hình khai triển của hình hộp chữ nhật có đáy là hình vuông cạnh $5\text{ cm}$, chiều cao $10\text{ cm}$. Bốn mặt xung quanh ghép liền nhau tạo thành một hình chữ nhật lớn có kích thước là
\choice
{$5\text{ cm}$ và $10\text{ cm}$}
{\True $20\text{ cm}$ và $10\text{ cm}$}
{$15\text{ cm}$ và $10\text{ cm}$}
{$25\text{ cm}$ và $10\text{ cm}$}
\loigiai{
Bốn mặt bên ghép liên tiếp tạo thành hình chữ nhật có:
- Chiều rộng là chiều cao hình hộp: $10\text{ cm}$.
- Chiều dài là chu vi đáy: $5 \times 4 = 20\text{ cm}$.
Vậy kích thước là $20\text{ cm}$ và $10\text{ cm}$.\\
\textbf{Đáp án của câu hỏi là B}
}
\end{ex}

\begin{ex} %[5H9V49]
Người ta dùng một tấm bìa hình chữ nhật có kích thước dài $30\text{ cm}$, rộng $20\text{ cm}$ để cắt thành hình khai triển của một hình lập phương cạnh $5\text{ cm}$. Diện tích phần bìa còn thừa ít nhất là
\choice
{$400\text{ cm}^2$}
{\True $450\text{ cm}^2$}
{$350\text{ cm}^2$}
{$500\text{ cm}^2$}
\loigiai{
Diện tích của tấm bìa hình chữ nhật ban đầu là:
$$30 \times 20 = 600\text{ (cm}^2\text{)}.$$
Diện tích toàn phần của hình lập phương gồm $6$ mặt vuông cạnh $5\text{ cm}$ là:
$$5 \times 5 \times 6 = 150\text{ (cm}^2\text{)}.$$
Diện tích bìa còn thừa là:
$$600 - 150 = 450\text{ (cm}^2\text{)}.$$
\textbf{Đáp án của câu hỏi là B}
}
\end{ex}

\begin{ex} %[5H9V49]
Cho một ống bơ sữa đặc có dạng hình trụ, đường kính đáy $7\text{ cm}$ và chiều cao $8\text{ cm}$. Người ta dán một tờ giấy nhãn quanh thân ống sữa (không kể nắp và đáy). Lấy $\pi = \dfrac{22}{7}$, diện tích nhãn giấy đó là
\choice
{$154\text{ cm}^2$}
{\True $176\text{ cm}^2$}
{$198\text{ cm}^2$}
{$210\text{ cm}^2$}
\loigiai{
Chu vi đáy của ống bơ là:
$$7 \times \dfrac{22}{7} = 22\text{ (cm)}.$$
Tờ nhãn bọc quanh thân ống có dạng hình chữ nhật với chiều dài $22\text{ cm}$ và chiều rộng $8\text{ cm}$.\\
Diện tích tờ nhãn là:
$$22 \times 8 = 176\text{ (cm}^2\text{)}.$$
\textbf{Đáp án của câu hỏi là B}
}
\end{ex}

\begin{ex} %[5H9V49]
Trong các hình khai triển của hình lập phương có dạng mạng lưới gồm $6$ hình vuông, khẳng định nào sau đây là \textbf{sai}?
\choice
{Không có hình khai triển nào có $4$ hình vuông xếp thành một góc vuông $2 \times 2$}
{\True Luôn có đúng $4$ mặt xung quanh nằm trên một hàng thẳng}
{Một hình lập phương có tất cả $11$ dạng hình khai triển phẳng khác nhau}
{Hai mặt đối diện của hình lập phương không bao giờ có chung cạnh trong hình khai triển}
\loigiai{
Khẳng định B sai vì ngoài dạng $1-4-1$ (có $4$ mặt liên tiếp trên một hàng thẳng), còn có các dạng như $2-3-1$, $2-2-2$, $3-3$ không có $4$ mặt liên tiếp trên một hàng thẳng.\\
\textbf{Đáp án của câu hỏi là B}
}
\end{ex}

\begin{ex} %[5H9V49]
Một hình khai triển gồm $6$ hình vuông bằng nhau có tổng chu vi của các đường viền bao quanh bên ngoài là $70\text{ cm}$. Biết hình khai triển có dạng chữ thập ($1-4-1$) gồm $14$ đoạn cạnh của hình vuông ở đường bao ngoài. Thể tích của hình lập phương gấp được là
\choice
{$64\text{ cm}^3$}
{\True $125\text{ cm}^3$}
{$216\text{ cm}^3$}
{$343\text{ cm}^3$}
\loigiai{
Đường bao quanh hình khai triển chữ thập gồm $14$ cạnh của các hình vuông nhỏ.\\
Độ dài mỗi cạnh hình vuông là:
$$70 : 14 = 5\text{ (cm)}.$$
Thể tích hình lập phương gấp được là:
$$5 \times 5 \times 5 = 125\text{ (cm}^3\text{)}.$$
\textbf{Đáp án của câu hỏi là B}
}
\end{ex}

\begin{ex} %[5H9V49]
Cắt các góc của một tấm bìa hình vuông cạnh $20\text{ cm}$ thành $4$ hình vuông nhỏ cạnh $4\text{ cm}$ rồi gấp các mép lên để được một chiếc hộp không nắp dạng hình hộp chữ nhật. Thể tích chiếc hộp đó là
\choice
{$640\text{ cm}^3$}
{\True $576\text{ cm}^3$}
{$720\text{ cm}^3$}
{$512\text{ cm}^3$}
\loigiai{
Sau khi cắt $4$ góc, đáy của chiếc hộp là hình vuông có cạnh:
$$20 - 4 \times 2 = 12\text{ (cm)}.$$
Chiều cao của chiếc hộp là $4\text{ cm}$.\\
Thể tích chiếc hộp là:
$$12 \times 12 \times 4 = 576\text{ (cm}^3\text{)}.$$
\textbf{Đáp án của câu hỏi là B}
}
\end{ex}

\textbf{Phần 2. Trắc nghiệm đúng sai}
\setcounterpageref{ex}{0}

\begin{ex} %[5H9V49]
Cho hình khai triển của một hình hộp chữ nhật gồm $6$ mặt có các kích thước đôi một là: hai mặt $6\text{ cm} \times 4\text{ cm}$, hai mặt $6\text{ cm} \times 3\text{ cm}$, hai mặt $4\text{ cm} \times 3\text{ cm}$. Xét tính đúng/sai của các nhận định sau:
\choiceTF
{\True Ba kích thước của hình hộp chữ nhật đó là $6\text{ cm}$, $4\text{ cm}$ và $3\text{ cm}$}
{\True Tổng diện tích của $6$ mặt trong hình khai triển là $108\text{ cm}^2$}
{Hình hộp chữ nhật này có diện tích xung quanh lớn nhất bằng $72\text{ cm}^2$}
{\True Thể tích của hình hộp chữ nhật gấp được là $72\text{ cm}^3$}
\loigiai{
\begin{itemchoice}
\item Đúng. Ba kích thước của hình hộp là chiều dài $6\text{ cm}$, chiều rộng $4\text{ cm}$, chiều cao $3\text{ cm}$ (hoặc hoán vị).
\item Đúng. Diện tích toàn phần: $(6 \times 4 + 6 \times 3 + 4 \times 3) \times 2 = (24 + 18 + 12) \times 2 = 108\text{ (cm}^2\text{)}$.
\item Sai. Nếu chọn chiều cao là $6\text{ cm}$ thì $S_{xq} = (4 + 3) \times 2 \times 6 = 84\text{ cm}^2 > 72\text{ cm}^2$.
\item Đúng. Thể tích hình hộp là $6 \times 4 \times 3 = 72\text{ (cm}^3\text{)}$.
\end{itemchoice}
}
\end{ex}

\begin{ex} %[5H9V49]
Một hình khai triển của một lon nước ngọt hình trụ gồm một hình chữ nhật kích thước $15{,}7\text{ cm} \times 12\text{ cm}$ và hai hình tròn bằng nhau. Biết chiều dài hình chữ nhật là chu vi đường tròn đáy (lấy $\pi = 3{,}14$). Xét tính đúng/sai của các nhận định sau:
\choiceTF
{\True Chiều cao của lon nước ngọt là $12\text{ cm}$}
{\True Đường kính của đáy lon là $5\text{ cm}$}
{Bán kính đáy của lon nước ngọt là $3\text{ cm}$}
{\True Diện tích xung quanh của lon nước ngọt là $188{,}4\text{ cm}^2$}
\loigiai{
\begin{itemchoice}
\item Đúng. Chiều rộng của hình chữ nhật là chiều cao của hình trụ: $h = 12\text{ cm}$.
\item Đúng. Đường kính đáy là: $d = 15{,}7 : 3{,}14 = 5\text{ (cm)}$.
\item Sai. Bán kính đáy là $r = 5 : 2 = 2{,}5\text{ cm}$.
\item Đúng. Diện tích xung quanh bằng diện tích hình chữ nhật: $15{,}7 \times 12 = 188{,}4\text{ (cm}^2\text{)}$.
\end{itemchoice}
}
\end{ex}

\textbf{Phần 3. Câu trả lời ngắn}
\setcounterpageref{ex}{0}

\begin{ex} %[5H9V49]
Một hình khai triển của hình lập phương có tổng chu vi tất cả $6$ hình vuông rời rạc là $96\text{ cm}$. Khi gấp lại thành hình lập phương thì thể tích của khối lập phương đó là bao nhiêu xăng-ti-mét khối?
\shortans{64}
\loigiai{
Chu vi của mỗi hình vuông là:
$$96 : 6 = 16\text{ (cm)}.$$
Cạnh của hình vuông (cũng là cạnh hình lập phương) là:
$$16 : 4 = 4\text{ (cm)}.$$
Thể tích hình lập phương là:
$$4 \times 4 \times 4 = 64\text{ (cm}^3\text{)}.$$
}
\end{ex}

\begin{ex} %[5H9V49]
Một hộp quà hình lập phương có cạnh $15\text{ cm}$. Người ta muốn cắt tấm giấy gói quà có diện tích đúng bằng diện tích toàn phần của hộp quà đó. Hỏi diện tích giấy cần dùng là bao nhiêu xăng-ti-mét vuông?
\shortans{1350}
\loigiai{
Diện tích một mặt của hộp quà là:
$$15 \times 15 = 225\text{ (cm}^2\text{)}.$$
Diện tích giấy cần dùng (diện tích toàn phần) là:
$$225 \times 6 = 1\,350\text{ (cm}^2\text{)}.$$
}
\end{ex}

\begin{ex} %[5H9V49]
Một hình khai triển của hình trụ gồm một hình chữ nhật có diện tích $200\text{ cm}^2$ và hai hình tròn đáy. Biết chiều rộng của hình chữ nhật bằng $8\text{ cm}$. Chu vi của mỗi đáy hình trụ là bao nhiêu xăng-ti-mét?
\shortans{25}
\loigiai{
Chiều dài của hình chữ nhật chính là chu vi đáy hình trụ:\\
$$200 : 8 = 25\text{ (cm)}.$$
}
\end{ex}

\begin{ex} %[5H9V49]
Bác Thợ mộc có một khối gỗ hình hộp chữ nhật dài $8\text{ dm}$, rộng $6\text{ dm}$, cao $5\text{ dm}$. Tổng diện tích tất cả các mặt bên ngoài của khối gỗ đó là bao nhiêu đề-xi-mét vuông?
\shortans{236}
\loigiai{
Diện tích xung quanh của khối gỗ là:
$$(8 + 6) \times 2 \times 5 = 140\text{ (dm}^2\text{)}.$$
Diện tích hai đáy của khối gỗ là:
$$8 \times 6 \times 2 = 96\text{ (dm}^2\text{)}.$$
Tổng diện tích tất cả các mặt là:
$$140 + 96 = 236\text{ (dm}^2\text{)}.$$
}
\end{ex}

\begin{ex} %[5H9V49]
Một hình lập phương có diện tích toàn phần là $600\text{ cm}^2$. Thể tích của hình lập phương đó là bao nhiêu xăng-ti-mét khối?
\shortans{1000}
\loigiai{
Diện tích một mặt của hình lập phương là:
$$600 : 6 = 100\text{ (cm}^2\text{)}.$$
Vì $10 \times 10 = 100$ nên cạnh hình lập phương là $10\text{ cm}$.\\
Thể tích hình lập phương là:
$$10 \times 10 \times 10 = 1\,000\text{ (cm}^3\text{)}.$$
}
\end{ex}

\begin{ex} %[5H9V49]
Từ một mảnh bìa chữ nhật có chiều dài $24\text{ cm}$ và chiều rộng $14\text{ cm}$, người ta cắt bỏ ở $4$ góc bốn hình vuông cạnh $2\text{ cm}$ rồi gấp lên thành chiếc hộp không nắp. Diện tích xung quanh của chiếc hộp đó là bao nhiêu xăng-ti-mét vuông?
\shortans{120}
\loigiai{
Chiều dài đáy hộp là: $24 - 2 \times 2 = 20\text{ (cm)}$.\\
Chiều rộng đáy hộp là: $14 - 2 \times 2 = 10\text{ (cm)}$.\\
Chiều cao hộp là $2\text{ cm}$.\\
Diện tích xung quanh của hộp là:
$$(20 + 10) \times 2 \times 2 = 120\text{ (cm}^2\text{)}.$$
}
\end{ex}
